\section{Grundlagen und Analyse}
\setauthor{\secondauthor}

\subsection{Marktanalyse}
In der Analyse bestehender digitaler Angebote liegt der Fokus auf deren Eignung für die spezifischen Anforderungen der Matura-Vorbereitung. Dabei spielen die authentische Prüfungssimulation, die Bewältigung des enormen Stoffumfangs sowie die fachliche Komplexität der Aufgabenstellungen eine entscheidende Rolle.

\subsubsection{StudySmarter}
StudySmarter wurde ursprünglich als Tool für das \textit{Microlearning} und zum Auswendiglernen von isoliertem Faktenwissen konzipiert. Diese historische Ausrichtung auf eine „häppchenweise“ Wissensvermittlung macht die Plattform zwar effizient für Vokabeltraining oder kurze Definitionen, führt jedoch bei der komplexen Matura-Vorbereitung an deutliche Grenzen.

Hinsichtlich der \textbf{Lern- und Quizfunktionen} setzt StudySmarter primär auf das Prinzip der digitalen Karteikarten. Benutzer können Sätze von Karten durcharbeiten, wobei ein Algorithmus (Spaced Repetition) die Abstände der Wiederholungen steuert. Die Quizzes bestehen meist aus einfachen Zuordnungen oder Single-Choice-Fragen, die auf ein schnelles Erfolgserlebnis abzielen. Eine gezielte Auswahl spezifischer Themenpools, wie sie für die systematische Vorbereitung auf einzelne Matura-Kapitel nötig wäre, ist oft nur durch das manuelle Sortieren hunderter Einzelkarten möglich, was bei dem großen Stoffumfang der Matura sehr zeitintensiv ist.

Im Bereich des \textbf{Prüfungsmodus} zeigt sich ein deutliches Defizit: Die Quiz-Funktionen liefern meist ein unmittelbares Feedback nach jeder einzelnen Frage. Zwar bietet StudySmarter mittlerweile Funktionen an, die als „Prüfung“ betitelt werden, diese sind jedoch oft nur eine Aneinanderreihung der bekannten Karteikarten ohne echte Prüfungssimulation. Für die Matura ist hingegen die Simulation einer kompletten, mehrstündigen Klausur ohne Unterbrechung entscheidend. Das sofortige Anzeigen der Lösung stört den Konzentrationsfluss und verhindert das Training der mentalen Ausdauer sowie des Zeitmanagements, welche für eine reale Prüfungssituation essenziell sind.

\subsubsection{Serlo / Anton}
Diese Plattformen sind historisch im Bereich der Primar- und Sekundarstufe 1 verwurzelt, wo spielerisches Lernen (\textit{Gamification}) im Vordergrund steht. Das \textbf{Lernkonzept} basiert hier auf kleinen Lerneinheiten, die mit Belohnungen wie Sternen oder Punkten verknüpft sind.

Die Komplexität der Aufgaben beschränkt sich meist auf einfache Multiple-Choice-Fragen oder Lückentexte, die per Drag-and-Drop gelöst werden. Während dies für die Festigung von Basiswissen in unteren Schulstufen ideal ist, verlangt die Matura – insbesondere an technischen Schulen – das aktive Erarbeiten und Nachvollziehen mehrschrittiger Lösungswege sowie Freitext-Antworten. Ein \textbf{Prüfungsmodus unter Zeitdruck} oder eine Simulation, die über das spielerische Sammeln von Punkten hinausgeht, ist hier kaum vorhanden. Die \textbf{Statistiken} dienen eher der Motivation als der harten datengestützten Analyse von Wissenslücken, da sie nicht präzise aufzeigen, an welcher Stelle eines komplexen Lösungsweges ein Schüler gescheitert ist.

\subsubsection{Zusammenfassender Vergleich und Alleinstellungsmerkmale}
Im Gegensatz zu diesen allgemein gehaltenen Apps ist das vorliegende Projekt gezielt auf die strukturelle Tiefe der Matura ausgelegt. Während herkömmliche Tools oft generische Inhalte und einfache „Richtig/Falsch“-Statistiken nutzen, setzt dieses System auf folgende Alleinstellungsmerkmale:

\begin{itemize}
    \item \textbf{Detaillierte Lösungswege:} Diese wurden von Schülern für Schüler aufbereitet, um komplexe Fehlerquellen verständlich zu machen und Erklärungen auf Augenhöhe zu bieten.
    \item \textbf{Gezielte Themenpools:} Eine flexible Auswahl spezifischer Themengebiete ermöglicht es, den Fokus exakt auf individuelle Schwachstellen zu legen, anstatt starren Lernpfaden folgen zu müssen.
    \item \textbf{Authentischer Prüfungsmodus:} Durch einstellbare Zeitlimits, freie Navigation zwischen den Fragen und verzögertes Feedback wird die reale Matura-Situation realitätsnah simuliert.
    \item \textbf{Übungsmodus mit direktem Feedback:} Dieser Modus ermöglicht es, sofortige Rückmeldung zu erhalten, um gezielt an Fehlern zu arbeiten, ohne den Druck einer kompletten Prüfungssimulation.
\end{itemize}

\subsection{Maturatraining digitalisieren}
\setauthor{\secondauthor}

Die Digitalisierung des Maturatrainings zielt darauf ab, fragmentierte Lernprozesse in eine strukturierte, interaktive Umgebung zu überführen. Ein zentrales Problem in der Vorbereitung war bisher die mangelnde Konsistenz der Unterlagen: Schüler:innen arbeiteten oft mit unvollständigen Mitschriften, die teils analog auf Papier und teils digital als unsortierte Fragmente vorlagen.

Das Partnerprojekt hat hierfür die notwendige Grundlage geschaffen, indem verstreute Mitschriften zentralisiert und in digitale Fragenkataloge überführt wurden. Da zu einer vollständigen Lernplattform ein Übungs- und Prüfungsmodus inkludiert sein sollte, konzentriert sich die vorliegende Arbeit auf die methodische Nutzbarmachung dieser Inhalte durch die Entwicklung eines Übungs- und Prüfungsmodus.